\begin{abstract}
Self-supervised pre-training works best with millions of images and the
compute to match, which puts visual foundation models out of reach for most
groups. We characterize how one widely used self-supervised method, Barlow
Twins, behaves when both data and compute are small. A 5.4M-parameter vision
transformer is pre-trained on subsets of Tiny-ImageNet (1k to 100k images);
each resulting network, its weights then frozen, is scored by training a
linear classifier on CIFAR-10 --- a downstream task the network never saw.
Checkpoints are selected with a validation score computed only on the
pre-training dataset, never on the downstream task. In the completed
single-seed campaign, pre-training pays at every set size: 61.4\% downstream
accuracy from only 1k unlabeled images and 71.5\% from 100k, against a
41.5\% floor for an untrained network. The central finding is that the
pre-training validation score and downstream accuracy disagree about
\emph{which checkpoint is best}: downstream-blind selection picks
checkpoints that score up to 4.7 points worse on a CIFAR-10 monitoring
metric than the best checkpoint in hindsight. The disagreement peaks at the
mid-range set sizes, where the measured curve even dips below its 1k value
--- the dip itself lies within the seed variance of an earlier campaign and
awaits confirmation; the divergence does not. On the smallest set,
validation accuracy decays 18\% from its peak with no sign of
representational collapse. At small scale, the rule for when to stop and
which checkpoint to keep shapes a data-efficiency curve alongside the data
itself.
\end{abstract}
